\section{Eigenvalues, Eigenvectors and Eigenspaces}
\subsection{Triangular and Diagonal Matrices/Linear Maps}

\begin{definition}[Triangular and Diagonal Matrices/Linear Maps]
    A matrix \(A \in \matrixset_{n \times n} \pqty{\FF}\) is \term{(upper) triangular} (respectively, \term{diagonal}) if \(A\) is of the form
    \[
        A = \pmqty{
            \lambda_1 & * & \cdots & * \\
             & \lambda_2 & \cdots & * \\
             & & \ddots & \vdots \\
            0 & & & \lambda_n
        },
    \]
    and respectively,
    \[
        A = \pmqty{
            \lambda_1 &  & & 0 \\
             & \lambda_2 & &  \\
             & & \ddots & \\
            0 & & & \lambda_n
        }.
    \]

    A matrix \(A \in \matrixset_{n \times n} \pqty{\FF}\) is \term{triangularisable} (respectively, \term{diagonalisable}) if it is similar to a triangular (respectively, diagonal) matrix.

    Let \(V\) be a finite dimensional vector space over \(\FF\), and let \(\alpha \in \linearmap\pqty{V, V}\). We say \(\alpha\) is \term{triangularisable} (respectively, \term{diagonalisable}) if there exists a basis \(\basis\) of \(V\) such that the matrix representation of \(\alpha\) with respect to \(\basis\) is triangular (respectively, diagonal).
\end{definition}

\begin{remark}
    Note a linear endomorphism \(\alpha\) is triangularisable (respectively, diagonalisable) if and only if its matrix representation with respect to some (equivalently, any) basis is triangularisable (respectively, diagonalisable).
\end{remark}

It is easy to compute determinants of triangular matrices, while diagonal matrices are easier to analyse similarity class of a matrix.

Note that if there is a basis \(\basis = \set{\vb{v}_1, \ldots, \vb{v}_n}\) is a basis for \(V\), such that
\[
    \bqty{\alpha}_{\basis}^{\basis} = \pmqty{\dmat{\lambda_1, \ddots, \lambda_n}},
\]
then \(\alpha\pqty{\vb{v}_i} = \lambda_i \vb{v}_i\) for all \(i = 1, \ldots, n\).

\begin{definition}[Eigenvalues and Eigenvectors]
    Let \(V\) be a vector space of \(\FF\), and let \(\alpha \in \linearmap\pqty{V, V}\).
    \begin{itemize}
        \item \(\lambda \in \FF\) is an \term{eigenvalue} of \(\alpha\), if there exists a non-zero vector \(\vb{0} \neq \vb{v} \in V\) such that \(\alpha\pqty{\vb{v}} = \lambda \vb{v}\).
        \item Such a vector \(\vb{0} \neq \vb{v} \in V\) is called an \term{eigenvector} of \(\alpha\) corresponding to the eigenvalue \(\lambda\).
        \item The \term{eigenspace} of \(\alpha\) corresponding to the eigenvalue \(\lambda\) is the set
        \[
            V_{\lambda} = \set{\vb{v} \in V}{\alpha\pqty{\vb{v}} = \lambda \vb{v}}.
        \]

    \end{itemize}
\end{definition}

\begin{remark}
    \begin{itemize}
        \item Note that \(V_{\lambda}\) is a subspace of \(V\), since we have \(V_{\lambda} = \ker\pqty{\alpha - \lambda \identity_V}\).
        \item In the finite-dimensional case, \(\alpha\) is diagonalisable, if and only if there exists a basis of \(V\) consisting of eigenvectors of \(\alpha\).
        \item If \(\vb{v} \in V_{\mu}\) and \(\lambda \neq \mu\), then
        \[
            \pqty{\alpha - \lambda \identity_V}\pqty{\vb{v}} = (\mu - \lambda) \vb{v},
        \]
        or in other words, the map \(\pqty{\alpha - \lambda \identity_V}\) preserves the subspace \(V_{\mu}\), and the restriction
        \[
            \functiondomain{\rightbar{\pqty{\alpha - \lambda \identity_V}}_{V_{\mu}}}{V_{\mu}}{V_{\mu}}
        \]
        is invertible, with inverse
        \[
            v \mapsto \pqty{\mu - \lambda}\inverse v.
        \]
    \end{itemize}
\end{remark}

\begin{lemma}[Eigenvectors corresponding to Distinct Eigenvalues are Linearly Independent]
    \label{lem:eigenvectors-distinct-eigenvalues-linearly-independent}%
    If \(\lambda_1, \ldots, \lambda_n\) are distinct eigenvalues of \(\alpha\), with corresponding eigenvectors \(\vb{0} \neq \vb{v}_i \in V_{\lambda_i}\). Then the set \(\set{\vb{v}_1, \ldots, \vb{v}_n}\) is linearly independent.
\end{lemma}

\begin{proof}
    We prove by contradiction. Suppose it is not linearly independent. Let
    \[
        \sum_{i = 1}^{n} \mu_i \vb{v}_i = \vb{0},
    \]
    for \(\mu_i \in \FF\), not all zero. Without loss of generality, let \(\mu_1 \neq 0\). Let
    \[
        \beta = \prod_{i = 2}^{n} \pqty{\alpha - \lambda_i \identity_V}.
    \]

    Then we have
    \[
        \beta\pqty{\vb{v}_j} = \pqty{\prod_{i = 2}^{n} \pqty{\lambda_j - \lambda_i}} \vb{v}_j
    \]
    and this is non-zero if and only if \(j = 1\). Thus,
    \begin{align*}
        \vb{0} &= \beta\pqty{\sum_{i = 1}^{n} \mu_i \vb{v}_i} \\
        &= \mu_1 \pqty{\prod_{i = 2}^{n} \pqty{\lambda_1 - \lambda_i}} \vb{v}_1, \\
        &\neq \vb{0},
    \end{align*}
    which is a contradiction.
\end{proof}

\begin{corollary}[Finitely Many Eigenvalues]
    If \(V\) is a finite-dimensional vector space, then \(\alpha \in \linearmap\pqty{V, V}\) has finitely many eigenvalues, and in particular it has at most \(\dim V\) distinct eigenvalues.
\end{corollary}

\begin{example}[Real Sequences and Infinitely Many Eigenvalues]
    Let \(V = \RR^{\NN}\), and define the linear map
    \[
        \function{\alpha}{V}{V}{\pqty{x_n}_{n \in \NN}}{\pqty{a_{n + 1}}_{n \in \NN}}.
    \]

    Then for every \(\lambda \in \RR\), the sequence \(\pqty{\lambda^n}_{n \in \NN}\) is an eigenvector of \(\alpha\) corresponding to the eigenvalue \(\lambda\). Thus, \(\alpha\) has infinitely many eigenvalues.
\end{example}

\begin{proposition}[Eigenspaces Generate Direct Sums]
    \label{prop:eigenspaces-generate-direct-sums}%
    Let \(V\) be a finite-dimensional vector space, and let \(\alpha \in \linearmap\pqty{V, V}\). Suppose \(\lambda_1, \ldots, \lambda_k \in \FF\) are all the eigenvalues of \(\alpha\).

    \begin{enumerate}[label=(\roman*)]
        \item \term{Eigenspaces always generate direct sums.} We have
        \[
            \spn\pqty{V_{\lambda_1} \cup \cdots \cup V_{\lambda_k}} = V_{\lambda_1} \oplus \cdots \oplus V_{\lambda_k}.
        \]

        \item If \(V\) is finite-dimensional, then \(\alpha\) is diagonalisable if and only if
        \[
            V = V_{\lambda_1} \oplus \cdots \oplus V_{\lambda_k}.
        \]
    \end{enumerate}
\end{proposition}

\begin{proof}
    \begin{enumerate}[label=(\roman*)]
        \item If this is not true, then there is some \(\vb{v}_i \in V_{\lambda_i}\) such that
        \[
            \sum_{i = 1}^{k} \vb{v}_i = \vb{0},
        \]
        and not all \(\vb{v}_i\) are zero. This contradicts \autoref{lem:eigenvectors-distinct-eigenvalues-linearly-independent}.

        \item Suppose \(\alpha\) is diagonalisable. Then there is a basis of eigenvectors of \(\alpha\). These eigenvectors span \(V\), and lie in \(V_{\lambda_1} \cup \cdots \cup V_{\lambda_k}\).
        
        Conversely, if \(V = V_{\lambda_1} \oplus \cdots \oplus V_{\lambda_k}\), and let \(\basis_i\) be a basis of \(V_{\lambda_i}\), with \(\abs{\basis_i} = \dim V_{\lambda_i} = n_i\). Then \(\basis = \bigcup_{i = 1}^{k} \basis_i\) is a basis of \(V\), and we have
        \[
            \bqty{\alpha}_{\basis}^{\basis} = \pmqty{\dmat{\lambda_1 I_{n_1}, \lambda_2 I_{n_2}, \ddots,\lambda_k I_{n_k}}}
        \]
        is diagonal, and thus \(\alpha\) is diagonalisable.
    \end{enumerate}
\end{proof}

\subsection{Characteristic Polynomial and Eigenvalues}

Recall that if \(V\) is finite-dimensional, and \(\beta \in \linearmap\pqty{V, V}\), then \(\det \beta = 0\) if and only if \(\ker \pqty{\beta} \neq \set{\vb{0}}\). Therefore, \(\lambda \in \FF\) is an eigenvalue of \(\alpha \in \linearmap\pqty{V, V}\) if and only if \(\det\pqty{\alpha - \lambda \identity_V} = 0\).

\begin{definition}[Characteristic Polynomial]
    Let \(A \in \matrixset_{n \times n} \pqty{\FF}\). The \term{characteristic polynomial} of \(A\) is the polynomial
    \[
        \characteristic_{A} \pqty{t} = \det\pqty{t I_n - A} \in \FF\bqty{t}.
    \]
    
    Similarly, for a finite-dimensional vector space \(V\) and \(\alpha \in \linearmap\pqty{V, V}\), the characteristic polynomial of \(\alpha\) is defined as
    \[
        \characteristic_{\alpha} \pqty{t} = \chi_{\bqty{\alpha}_{\basis}^{\basis}} \pqty{t},
    \]
    for some (equivalently, any) basis \(\basis\) of \(V\).
\end{definition}

\begin{remark}
    \begin{enumerate}[label=(\roman*)]
        \item The characteristic polynomial \(\characteristic_{A}\) is monic of degree \(n\).
        \item Note that
        \[
            \characteristic_{P A P\inverse} \pqty{t} = \det\pqty{t I_n - P A P\inverse} = \det\pqty{P\pqty{t I_n - A}P\inverse} = \det\pqty{t I_n - A} = \characteristic_{A}\pqty{t},
        \]
        so similar matrices have the same characteristic polynomial. In particular, \(\characteristic_{\alpha}\) is well-defined.
        \item By the Leibniz rule, if the characteristic polynomial is given by
        \[
            \characteristic_{\alpha} \pqty{t} = t^n + c_{n - 1} t^{n - 1} + \cdots + c_1 t + c_0,
        \]
        then \(c_{n - 1} = -\tr \alpha\) and \(c_0 = (-1)^n \det \alpha\).
    \end{enumerate}
\end{remark}

Then our previous observation of eigenvalues and characteristic polynomial can be summarised as follows.

\begin{lemma}[Eigenvalues are Roots of Characteristic Polynomial]
    Let \(V\) be a finite-dimensional vector space, and let \(\alpha \in \linearmap\pqty{V, V}\). The set of eigenvalues of \(\alpha\) is precisely the set of roots of \(\characteristic_{\alpha}\).
\end{lemma}

\subsection{Triangularisability and Diagonalisability}

\begin{theorem}[Triangularisability of Complex Matrices]
    \label{thm:triangularisability-complex-matrices}%
    Every square, complex matrix is triangularisable.
\end{theorem}

\begin{proposition}[Triangularisability and Characteristic Polynomial]
    \label{prop:triangularisability-characteristic-polynomial}%
    For any \(A \in \matrixset_{n \times n} \pqty{\FF}\), \(A\) is triangularisable if and only if \(\characteristic_{A}\) splits into linear factors over \(\FF\).
\end{proposition}

Note that \autoref{thm:triangularisability-complex-matrices} follows from \autoref{prop:triangularisability-characteristic-polynomial}, by \autoref{thm:fundamental-theorem-of-algebra} (Fundamental Theorem of Algebra).

\begin{example}[Real Matrices and Non-Triangularisability]
    Note that the matrix
    \[
        A = \pmqty{\cos \theta & -\sin \theta \\ \sin \theta & \cos \theta} \in \matrixset_{2 \times 2} \pqty{\RR},
    \]
    with \(\theta \in \oointv{0}{\pi}\). Then the characteristic polynomial is
    \[
        \characteristic_{A} \pqty{t} = t^2 - 2 \cos \theta t + 1,
    \]
    which has discriminant \(\Delta = 4 \cos^2 \theta - 4 < 0\). Thus, \(\characteristic_{A}\) does not split into linear factors over \(\RR\), and thus \(A\) is not triangularisable over \(\RR\).
\end{example}

\begin{proof}[\proofname\ of \autoref{prop:triangularisability-characteristic-polynomial}]
    Let \(\functiondomain{\alpha}{\FF^n}{\FF^n}\) be given by \(\alpha\pqty{\vb{v}} = A\vb{v}\), so that \(\bqty{\alpha}_{E_n}^{E_n} = A\).

    Suppose \(A\) is triangularisable. Then A is similar to
    \[
        A' = \pmqty{
            \lambda_1 & * & \cdots & * \\
             & \lambda_2 & \cdots & * \\
             & & \ddots & \vdots \\
            0 & & & \lambda_n
        }.
    \]

    Then we have
    \[
        \characteristic_{A} \pqty{t} = \characteristic_{A'} \pqty{t} = \det\pqty{t I_n - A'} = \prod_{i = 1}^{n} \pqty{t - \lambda_i}.
    \]

    Conversely, suppose \(\characteristic_{A}\) splits into linear factors over \(\FF\). Then in particular, \(\characteristic_{A}\) has a root \(\lambda \in \FF\). Then \(\alpha\) has an eigenvector \(\vb{0} \neq \vb{v} \in V_{\lambda}\) corresponding to \(\lambda\).

    Extend \(\set{\vb{v}_1}\) to a basis \(\basis = \set{\vb{v}_1, \ldots, \vb{v}_n}\) of \(\FF^n\). Then by change of basis matrix, \(A\) is similar to
    \[
        \bqty{\alpha}_{\basis}^{\basis} = \pmqty{
            \lambda & * & \cdots & * \\
            0 &  &  &  \\
            \vdots &  & C &  \\
            0 &  &  &
        },
    \]
    so the characteristic polynomial is
    \[
        \characteristic_{A} \pqty{t} = \pqty{t - \lambda} \characteristic_C \pqty{t},
    \]
    so \(\characteristic_C\) splits into linear factors over \(\FF\). By induction on the size of the matrices, we can conclude that \(C\) is triangularisable, with \(P C P\inverse\) being upper triangular. Let
    \[
        \overline{P} = \pmqty{1 & 0 \\ 0 & P},
    \]
    then \(A\) is similar to
    \[
        \overline{P} \bqty{\alpha}_{\basis}^{\basis} \overline{P}\inverse = \pmqty{
            \lambda & * & \cdots & * \\
            0 &  &  &  \\
            \vdots &  & P C P\inverse &  \\
            0 &  &  &
        },
    \]
    which is an upper triangular matrix. Thus, \(A\) is triangularisable.
\end{proof}

\subsection{Minimal Polynomials and the Cayley--Hamilton Theorem}

\begin{definition}[Polynomial on Linear Maps and Matrices]
    Let \(V\) be a finite-dimensional vector space over \(\alpha\), and let \(\alpha \in \linearmap\pqty{V, V}\). Let \(p\) be a polynomial
    \[
        p\pqty{t} = \sum_{i = 0}^{n} a_i t^i \in \FF\bqty{t}.
    \]

    We denote
    \[
        p\pqty{\alpha} = \sum_{i = 0}^{n} a_i \alpha^i \in \linearmap\pqty{V, V},
    \]
    where \(\alpha^0 = \identity_V\).

    Similarly, for \(A \in \matrixset_{n \times n} \pqty{\FF}\), we denote
    \[
        p\pqty{A} = \sum_{i = 0}^{n} a_i A^i \in \matrixset_{n \times n} \pqty{\FF},
    \]
    where \(A^0 = I_n\).
\end{definition}

\begin{remarks}
    Let \(V\) be finite-dimensional.
    \begin{itemize}
        \item For a basis \(\basis\) of \(V\), we have
        \[
            \bqty{p\pqty{\alpha}}_{\basis}^{\basis} = p\pqty{\bqty{\alpha}_{\basis}^{\basis}}.
        \]

        \item If \(\dim V = n\), then \(\dim \linearmap\pqty{V, V} = n^2\), so the set
        \[
            \set{\identity_V, \alpha, \ldots, \alpha^{n^2}} \subseteq \linearmap\pqty{V, V}
        \]
        has \(n^2 + 1\) elements, and thus is linearly dependent, so they satisfy a non-trivial linear combination
        \[
            0 = \sum_{i = 0}^{n^2} a_i \alpha^i = p\pqty{\alpha},
        \]
        where the polynomial is given by
        \[
            p\pqty{t} = \sum_{i = 0}^{n^2} a_i t^i \in \FF\bqty{t}.
        \]
    \end{itemize}
\end{remarks}

Denote
\[
    I_\alpha = \set{p \in \FF\bqty{t}}{p\pqty{\alpha} = 0} \subseteq \FF\bqty{t},
\]
then the remark above says that \(I_\alpha\) is non-empty, and it has an element of at most degree \(n^2\).

\begin{definition}[Minimal Polynomial]
    Let \(V\) be a finite-dimensional vector space over \(\FF\), and let \(\alpha \in \linearmap\pqty{V, V}\). The \term{minimal polynomial} of \(\alpha\) is the unique (non-zero) monic polynomial \(m_\alpha \in \FF\bqty{t}\) of minimal degree such that
    \[
        m_\alpha\pqty{\alpha} = 0.
    \]
\end{definition}

\begin{lemma}[Minimal Polynomial divides any Polynomial]
    \label{lem:minimal-polynomial-divides-any-polynomial}%
    Let \(V\) be a finite-dimensional vector space over \(\FF\), and let \(\alpha \in \linearmap\pqty{V, V}\). Then the minimal polynomial \(m_\alpha\) divides any polynomial \(f \in I_\alpha\).
\end{lemma}

\begin{proof}
    Suppose not, then we apply the Euclidean algorithm to \(m_\alpha\) and \(f\) to obtain
    \[
        f = m_\alpha q + r,
    \]
    where \(q, r \in \FF\bqty{t}\) and \(\deg r < \deg m_\alpha\) with \(r \neq 0\). Then we have
    \[
        0 = f\pqty{\alpha} = m_\alpha\pqty{\alpha} q\pqty{\alpha} + r\pqty{\alpha} = r\pqty{\alpha},
    \]
    so \(r \in I_\alpha\). So \(r \neq 0\) is a polynomial of smaller degree than \(m_\alpha\) such that \(r\pqty{\alpha} = 0\), which contradicts the minimality of \(m_\alpha\).
\end{proof}

\begin{proposition}[Uniqueness of Minimal Polynomial]
    Let \(V\) be a finite-dimensional vector space over \(\FF\), and let \(\alpha \in \linearmap\pqty{V, V}\). Then the minimal polynomial \(m_\alpha\) is unique.
\end{proposition}

\begin{proof}
    Suppose \(0 \neq \bar{m}_{\alpha} \in I_{\alpha}\) of minimal degree and is monic. By \autoref{lem:minimal-polynomial-divides-any-polynomial}, we have \(m_\alpha \divides \bar{m}_\alpha\), i.e.
    \[
        \bar{m}_\alpha = m_\alpha q
    \]
    for some \(q \in \FF\bqty{t}\). But \(\deg \pqty{m_\alpha} = \deg \pqty{\bar{m}_\alpha}\) by minimality, so \(q\) is a non-zero constant. Since both \(m_\alpha\) and \(\bar{m}_\alpha\) are monic, we have \(q = 1\), so \(m_\alpha = \bar{m}_\alpha\).
\end{proof}

\begin{theorem}[Cayley--Hamilton Theorem]
    \label{thm:cayley-hamilton-theorem}%
    Let \(V\) be a finite-dimensional vector space over \(\FF\), and let \(\alpha \in \linearmap\pqty{V, V}\). Then \(\alpha\) satisfies its own characteristic polynomial, i.e.
    \[
        \characteristic_{\alpha}\pqty{\alpha} = 0.
    \]
\end{theorem}

\begin{corollary}[Minimal Polynomial divides Characteristic Polynomial]
    Let \(V\) be a finite-dimensional vector space over \(\FF\), and let \(\alpha \in \linearmap\pqty{V, V}\). Then the minimal polynomial \(m_\alpha\) divides the characteristic polynomial \(\characteristic_{\alpha}\).
\end{corollary}

\begin{proof}[\proofname\ of \autoref{thm:cayley-hamilton-theorem} over \(\CC\)]
    By \autoref{thm:triangularisability-complex-matrices}, there is a basis
    \(\basis = \set{v_1, v_2, \ldots, v_n}\) of \(V\) such that
    \[
        \bqty{\alpha}_{\basis}^{\basis} = \pmqty{
            \lambda_1 & * & \cdots & * \\
             & \lambda_2 & \cdots & * \\
             & & \ddots & \vdots \\
            0 & & & \lambda_n
        },
    \]
    so we have
    \[
        \characteristic_{\alpha}\pqty{t} = \prod_{i = 1}^{n} \pqty{t - \lambda_i}.
    \]

    Define the subspace
    \[
        U_j = \spn\set{v_1, \ldots, v_j} \subspace V,
    \]
    so that \(U_0 = \set{\vb{0}}\) and \(U_n = V\). Then for each \(1 \leq j \leq n\), we have
    \[
        \pqty{\alpha - \lambda_j \identity_V}\pqty{U_j} \subspace U_{j - 1},
    \]
    so we have
    \[
        \characteristic_{\alpha} \pqty{\alpha} \pqty{V} = \pqty{\alpha - \lambda_1 \identity_V} \circ \cdots \circ \pqty{\alpha - \lambda_n \identity_V} \pqty{U_n} = \set{\vb{0}}.
    \]

    This shows that \(\characteristic_{\alpha}\pqty{\alpha} = 0\).
\end{proof}

\begin{remark}
    There is an alternative proof valid over any field \(\FF\) \textbf{\textcolor{red}{See Handout}}
\end{remark}

\subsection{Multiplicites and Diagonalisability}

\begin{definition}[Algebraic and Geometric Multiplicities]
    Let \(V\) be a finite-dimensional vector space over \(\FF\), and let \(\alpha \in \linearmap\pqty{V, V}\). Let \(\lambda \in \FF\) be an eigenvalue of \(\alpha\).
    
    The \term{algebraic multiplicity}, denoted \(a_{\lambda}\), of \(\lambda\) is the multiplicity of \(\lambda\) as a root to the characteristic polynomial \(\characteristic_{\alpha}\), i.e., the largest integer \(k\) such that \(\pqty{t - \lambda}^k\) divides \(\characteristic_{\alpha}\).

    The \term{geometric multiplicity}, denoted \(g_{\lambda}\), of \(\lambda\) is the dimension of the eigenspace \(V_{\lambda}\).
\end{definition}

Note that \(g_{\lambda}, a_{\lambda} \geq 1\).

\begin{proposition}[Algebraic and Geometric Multiplicities]
    \label{prop:algebraic-geometric-multiplicities}%
    Let \(V\) be a finite-dimensional vector space over \(\FF\), and let \(\alpha \in \linearmap\pqty{V, V}\). Let \(\lambda \in \FF\) be an eigenvalue of \(\alpha\). Then
    \[
        g_{\lambda} \leq a_{\lambda}.
    \]
\end{proposition}

\begin{proof}
    Let \(\basis_\lambda\) be a basis for \(V_{\lambda}\), and we extend to a basis \(\basis \) of \(V\). Then we have the matrix
    \[
        \bqty{\alpha}_{\basis}^{\basis} = \pmqty{
            \lambda I_{g_\lambda} & * \\
            0 & C
        },
    \]  
    and hence
    \[
        \characteristic_{\alpha} \pqty{t} = \det\pqty{t I_n - \bqty{\alpha}_{\basis}^{\basis}} = \pqty{t - \lambda}^{g_\lambda} \characteristic_C \pqty{t},
    \]
    so at least \(g_\lambda\) factors of \(\pqty{t - \lambda}\) divide \(\characteristic_{\alpha}\), so \(g_\lambda \leq a_\lambda\).
\end{proof}

\begin{lemma}
    \label{lem:sum-of-multiplicities}%
    Let \(V\) be a finite-dimensional vector space over \(\FF\), and let \(\alpha \in \linearmap\pqty{V, V}\). Let the eigenvalues of \(\alpha\) be \(\lambda_1, \ldots, \lambda_k\). Then we have the following.

    \begin{enumerate}[label=(\roman*)]
        \item \textbf{Sum of Geometric Multiplicities.}
        \[
            \sum_{i = 1}^{k} g_{\lambda_i} \leq \dim V,
        \]
        with equality if and only if \(\alpha\) is diagonalisable.

        \item \textbf{Sum of Algebraic Multiplicities.}
        \[
            \sum_{i = 1}^{k} a_{\lambda_i} \leq \dim V.
        \]
        with equality if and only if \(\alpha\) is triangularisable.
    \end{enumerate}
\end{lemma}

\begin{proof}
    \begin{enumerate}[label=(\roman*)]
        \item By \autoref{prop:eigenspaces-generate-direct-sums}, we have
        \[
            V_{\lambda_1} \oplus \cdots \oplus V_{\lambda_k} \subspace V,
        \]
        with equality if and only if \(\alpha\) is diagonalisable. Taking dimensions, we have
        \[
            \sum_{i = 1}^{k} g_{\lambda_i} = \sum_{i = 1}^{k} \dim V_{\lambda_i} = \dim\pqty{V_{\lambda_1} \oplus \cdots \oplus V_{\lambda_k}} \leq \dim V,
        \]
        with equality if and only if \(\alpha\) is diagonalisable.

        \item Write
        \[
            \characteristic_\alpha \pqty{t} = g\pqty{t} \cdot \prod_{i = 1}^{k} \pqty{t - \lambda_i}^{a_{\lambda_i}},
        \]
        then
        \begin{align*}
            \dim\pqty{V} &= \deg\pqty{\characteristic_\alpha} \\
            &= \deg\pqty{g} + \sum_{i = 1}^{k} a_{\lambda_i}.
        \end{align*}

        Since every root of \(\characteristic_{\alpha}\) is an eigenvalue (one of the \(\lambda_i\)), we have \(\deg g = 0\) if and only if \(\characteristic_{\alpha}\) splits into linear factors over \(\FF\), which is equivalent to \(\alpha\) being triangularisable by \autoref{prop:triangularisability-characteristic-polynomial}.
    \end{enumerate}
\end{proof}

\begin{proposition}[Diagonalisability and Multiplicities of Complex Matrices]
    \label{prop:diagonalisability-multiplicities-complex-matrices}%
    Let \(V\) be a finite-dimensional vector space over \(\CC\), and let \(\alpha \in \linearmap\pqty{V, V}\). Then \(\alpha\) is diagonalisable if and only if \(g_{\lambda} = a_{\lambda}\) for all eigenvalues \(\lambda\) of \(\alpha\).
\end{proposition}

\begin{proof}
    By \autoref{thm:triangularisability-complex-matrices}, \(\alpha\) is triangularisable, so we have \(\sum_{i = 1}^{k} a_{\lambda_i} = \dim V\), by \autoref{lem:sum-of-multiplicities} (ii).

    Since \(g_{\lambda_i} \leq a_{\lambda_i}\) for all \(i\) by \autoref{prop:algebraic-geometric-multiplicities}, we have
    \[
        \sum_{i = 1}^{k} g_{\lambda_i} \leq \sum_{i = 1}^{k} a_{\lambda_i} = \dim V,
    \]
    with equality if and only if \(g_{\lambda_i} = a_{\lambda_i}\) for all \(i\). By \autoref{lem:sum-of-multiplicities} (i), we have \(\alpha\) is diagonalisable if and only if \(\sum_{i = 1}^{k} g_{\lambda_i} = \dim V\), which is equivalent to \(g_{\lambda_i} = a_{\lambda_i}\) for all \(i\).
\end{proof}

\begin{remark}
    For any polynomial \(p\pqty{t} = \sum_{i = 0}^{n} \mu_i t^i \in \FF\bqty{t}\), and \(\vb{v} \in V_{\lambda}\), we have
    \[
        p\pqty{\alpha} \pqty{\vb{v}} = \sum_{i = 1}^{n} \mu_i \alpha^i\pqty{\vb{v}} = \sum_{i = 1}^{n} \mu_i \lambda^i \vb{v} = p\pqty{\lambda} \vb{v}.
    \]
\end{remark}

\begin{lemma}[Minimal and Characteristic Polynomials have the Same Roots]
    For \(\lambda \in \FF\), \(\lambda\) is a root of \(m_{\alpha}\) if and only if \(\lambda\) is a root of \(\characteristic_{\alpha}\).
\end{lemma}

\begin{proof}
    If \(\lambda\) is a root of \(m_{\alpha}\), then the Cayley--Hamilton thorem implies that \(\lambda\) is a root of \(\characteristic_{\alpha}\).

    If \(\lambda\) is a root of \(\characteristic_{\alpha}\), then \(\lambda\) is an eigenvalue of \(\alpha\). Take \(\vb{0} \neq \vb{v} \in V_{\lambda}\), then we have
    \[
        m_{\alpha} \pqty{\alpha} \pqty{\vb{v}} = m_{\alpha} \pqty{\lambda} \vb{v},
    \]
    while \(m_{\alpha} \pqty{\alpha} = 0\), so \(m_{\alpha} \pqty{\lambda} = 0\) since \(\vb{v} \neq \vb{0}\), i.e. \(\lambda\) is a root of \(m_{\alpha}\).
\end{proof}

\begin{notation}
    Denote \(c_{\lambda}\) as the multiplicity of \(\lambda\) as a root of \(m_{\alpha}\). Observe \(1 \leq c_{\lambda} \leq a_{\lambda}\) for all eigenvalues \(\lambda\) of \(\alpha\).
\end{notation}

\begin{example}
    \begin{enumerate}[label=(\roman*)]
        \item Let \(A = \lambda I_n\). Then \(\characteristic_{A} \pqty{t} = \pqty{t - \lambda}^n\), \(V_{\lambda} = \FF^n\), \(m_{A}\pqty{t} = t - \lambda\). So \(g_{\lambda} = a_{\lambda} = n\) and \(c_{\lambda} = 1\).
        \item Let \(A\) be a Jordan block
        \[
            A = \pmqty{
                \lambda & 1 & & & 0 \\
                 & \lambda & 1 & & \\
                 &  & \ddots & \ddots & \\
                 & & & \lambda & 1 \\
                0 & & & & \lambda
            }.
        \]

        Then the characteristic polynomial is identical \(\characteristic_{A} \pqty{t} = \pqty{t - \lambda}^n\), however the eigenspace is now only \(V_{\lambda} = \spn\set{\vb{e}_1}\). Thus, \(g_{\lambda} = 1\), \(a_{\lambda} = n\). Further, note
        \[
            \pqty{A - \lambda I_n}^{n - 1} = \pmqty{
                0 & \cdots & 0 & 1 \\
                0 & \cdots & 0 & 0 \\
                \vdots & \ddots & \vdots & \vdots \\
                0 & \cdots & 0 & 0
            },
        \]
        which is non-zero. So we must have \(m_A\pqty{t} = \pqty{t - \lambda}^n\), so \(c_{\lambda} = n\).
    \end{enumerate}
\end{example}

\begin{theorem}[Diagonalisability and Factorisation of Minimal Polynomial]
    \label{thm:diagonalisability-factorisation-minimal-polynomial}%
    Let \(V\) be a finite-dimensional vector space over \(\FF\), and let \(\alpha \in \linearmap\pqty{V, V}\). If \(p\pqty{\alpha} = 0\) for some \(0 \neq p\pqty{t} \in \FF\bqty{t}\) which is a product of distinct linear factors, then \(\alpha\) is diagonalisable.
\end{theorem}

\begin{example}
    Take \(\FF = \CC\). Note every invertible matrix \(A \in \GL_n\pqty{\CC}\) of finite order is diagonalisable. Indeed, suppose \(A^d = I\). Set \(p\pqty{t} = t^d - 1 \in \CC\bqty{t}\). Then this splits up as linear factors
    \[
        p\pqty{t} = \prod_{i = 1}^{d} \pqty{t - \omega^i} \qfor \omega = \exp\pqty{\frac{2 \pi i}{d}}.
    \]

    So \(p\pqty{A} = 0\), and hence \(A\) is diagonalisable.
\end{example}

\begin{proof}[\proofname\ of \autoref{thm:diagonalisability-factorisation-minimal-polynomial}]
    We may assume that \(p\) is monic. Then write
    \[
        p\pqty{t} = \prod_{i = 1}^{k} \pqty{t - \mu_i}
    \]
    for some distinct \(\mu_i \in \FF\). Set polynomials
    \[
        p_i\pqty{t} = \prod_{j \neq i} \pqty{t - \mu_j}.
    \]  

    Then the \(p_i\)s are pairwise coprime (no common factors of degree at least one, since the \(\mu_i\)s are distinct). Appplying B\'{e}zout's lemma, there exist \(q_i\pqty{t} \in \FF\bqty{t}\), such that
    \[
        1 = \sum_{i = 1}^{k} p_i\pqty{t} q_i\pqty{t}.
    \]

    Take \(\pi_i = p_i\pqty{\alpha} \circ q_i\pqty{\alpha} \in \linearmap\pqty{V, V}\). Then
    \[
        \identity_V = \sum_{i = 1}^{k} \pi_i,
    \]
    so for each \(\vb{v} \in V\), we have
    \[
        \vb{v} = \sum_{i = 1}^{k} \pi_i \pqty{\vb{v}},
    \]
    and hence
    \[
        V = \img\pqty{\pi_1} + \cdots + \img\pqty{\pi_k}.
    \]

    Note that \(\pqty{t - \mu_i} p_i\pqty{t} = p\pqty{t}\), so in terms of endomorphisms, we have
    \[
        \pqty{\alpha - \mu_i \identity_V} \circ \pi_i = p\pqty{\alpha} \circ q_i \pqty{\alpha} = 0,
    \]
    so \(\img\pqty{\pi_i} \subspace \ker\pqty{\alpha - \mu_i \identity_V} = V_{\mu_i}\) Thus we have
    \begin{align*}
        V &= \img\pqty{\pi_1} + \cdots + \img\pqty{\pi_k} \\
        &\subspace V_{\mu_1} \oplus \cdots \oplus V_{\mu_k} \\
        &\subspace V.
    \end{align*}

    Hence, \(V = V_{\mu_1} \oplus \cdots \oplus V_{\mu_k}\), and thus \(\alpha\) is diagonalisable.
\end{proof}

So far, we summarise criterion for diagonalisability.

\begin{theorem}[Criterion for Diagonalisability]
    Let \(V\) be a finite-dimensional vector space over \(\FF\), and let \(\alpha \in \linearmap\pqty{V, V}\). Then the following are equivalent.
    \begin{enumerate}[label=(\roman*)]
        \item \(\alpha\) is diagonalisable.
        \item \(V\) has a basis of eigenvectors of \(\alpha\).
        \item There exists \(0 \neq p\pqty{t} \in \FF\bqty{t}\) which is a product of distinct linear factors such that \(p\pqty{\alpha} = 0\).
        \item The minimal polynomial \(m_\alpha\) splits into distinct linear factors.
        \item If \(\FF = \CC\), for all eigenvalues \(\lambda\) of \(\alpha\), we have \(g_{\lambda} = a_{\lambda}\).
    \end{enumerate}
\end{theorem}

\begin{proof}
    \begin{itemize}
        \item (i) \(\iff\) (ii) is by definition.
        \item (i) \(\iff\) (v) in \(\FF = \CC\) is by \autoref{prop:diagonalisability-multiplicities-complex-matrices}.
        \item (iii) \(\implies\) (iv) is by \autoref{lem:minimal-polynomial-divides-any-polynomial}.
        \item (iv) \(\implies\) (iii) is trivial by simply taking \(p = m_\alpha\).
        \item (iii) \(\implies\) (i) is by \autoref{thm:diagonalisability-factorisation-minimal-polynomial}.
        \item (i) \(\implies\) (iii) is follows. Let \(\basis\) be a basis for \(V\) such that \(\bqty{\alpha}_{\basis}^{\basis}\) is diagonal. Then rearranging \(\basis\), we may assume that
        \[
            \bqty{\alpha}_{\basis}^{\basis} = \pmqty{
                \lambda_1 I_{n_1} &  & & 0 \\
                 & \lambda_2 I_{n_2} & &  \\
                 & & \ddots & \\
                0 &  &  & \lambda_k I_{n_k}
            },
        \]
        where \(\lambda_i \in \FF\) are distinct. Then the polynomial
        \[
            p\pqty{t} = \prod_{i = 1}^{k} \pqty{t - \lambda_i}
        \]
        is a product of distinct linear factors, and note that
        \[
            \bqty{p\pqty{\alpha}}_{\basis}^{\basis} = p\pqty{\bqty{\alpha}_{\basis}^{\basis}} = O,
        \]
        so \(p\pqty{\alpha} = 0\).
    \end{itemize}
\end{proof}

\begin{definition}[Simultaneously Diagonalisable]
    Let \(V\) be a finite-dimensional vector space over \(\FF\), and let \(\alpha, \beta \in \linearmap\pqty{V, V}\). We say that \(\alpha\) and \(\beta\) are \term{simultaneously diagonalisable} if there exists a basis \(\basis\) of \(V\) such that both \(\bqty{\alpha}_{\basis}^{\basis}\) and \(\bqty{\beta}_{\basis}^{\basis}\) are diagonal.
\end{definition}

\begin{lemma}[Commuting Endomorphisms preserves Eigenspaces]
    \label{lem:commuting-endomorphisms-preserves-eigenspaces}%
    Let \(V\) be a finite-dimensional vector space over \(\FF\), and let \(\alpha, \beta \in \linearmap\pqty{V, V}\) such that \(\alpha \circ \beta = \beta \circ \alpha\). Let \(V_{\lambda}\) be an eigenspace of \(\alpha\) corresponding to the eigenvalue \(\lambda\). Then \(\beta\pqty{V_{\lambda}} \subspace V_{\lambda}\).
\end{lemma}

\begin{proof}
    Let \(\vb{v} \in V_{\lambda}\). Note that
    \begin{align*}
        \alpha\pqty{\beta\pqty{\vb{v}}} &= \beta\pqty{\alpha\pqty{\vb{v}}} \\
        &= \beta \pqty{\lambda \vb{v}} \\
        &= \lambda \beta\pqty{\vb{v}},
    \end{align*}
    so \(\beta\pqty{\vb{v}} \in V_{\lambda}\).
\end{proof}

\begin{theorem}[Simultaneous Diagonalisability and Commutativity]
    Suppose \(\alpha\) and \(\beta\) are diagonalisable. They are simultaneously diagonalisable if and only if they commute, i.e. \(\alpha \circ \beta = \beta \circ \alpha\).
\end{theorem}

\begin{proof}
    Clearly commutativity is necessary, since diagonal matrices commute.

    Suppose \(\alpha\) and \(\beta\) are diagonalisable and commute. Let
    \[
        V = V_{\lambda_1} \oplus \cdots \oplus V_{\lambda_k}
    \]
    for \(V_{\lambda_i}\) being the eigenspace of \(\alpha\) corresponding to the eigenvalue \(\lambda_i\). By \autoref{lem:commuting-endomorphisms-preserves-eigenspaces}, we have \(\beta\pqty{V_{\lambda_i}} \subspace V_{\lambda_i}\) for all \(i\). Let \(p\pqty{t} \in \FF\bqty{t}\) be a product of distinct linear factors such that \(p\pqty{\beta} = 0\). Then we may restrict \(\functiondomain{\beta}{V_{\lambda_i}}{V_{\lambda_i}}\), with \(p\pqty{\beta\vert_{V_{\lambda_i}}} = 0\). By \autoref{thm:diagonalisability-factorisation-minimal-polynomial}, we have \(\beta\vert_{V_{\lambda_i}}\) is diagonalisable, so there exists a basis \(\basis_i\) of \(V_{\lambda_i}\) such that \(\bqty{\beta}_{\basis_i}^{\basis_i}\) is diagonal. Then
    \[
        \basis = \bigcup_{i = 1}^{k} \basis_i
    \]
    is a basis for \(V\), and the elements of \(\basis\) are eigenvectors for both \(\alpha\) abd \(\beta\), so \(\alpha\) and \(\beta\) are simultaneously diagonalisable.
\end{proof}

\subsection{Jordan Normal Form}

\begin{definition}[Jordan Block]
    Let \(\lambda \in \CC\). The \(n \times n\) matrix
    \[
        J_n\pqty{\lambda} = \pmqty{
            \lambda & 1 & & & 0 \\
             & \lambda & 1 & & \\
             & & \ddots & \ddots & \\
             & & & \lambda & 1 \\
            0 & & & & \lambda
        }
    \]
    is called a \term{Jordan block} of size \(n\) corresponding to the eigenvalue \(\lambda\). We allow for \(n = 1\) in which case \(J_1\pqty{\lambda} = \pmqty{\lambda}\).
\end{definition}

Recall that the characteristic polynomial of \(J_n\pqty{\lambda}\) is \(\characteristic_{J_n\pqty{\lambda}} \pqty{t} = \pqty{t - \lambda}^n\), which is identical to its minimal polynomial. Its only eigenspace is \(V_{\lambda} = \spn\set{\vb{e}_1}\).

\begin{definition}[Jordan Normal Form]
    Let \(A \in \matrixset_{n \times n} \pqty{\CC}\). We say that \(A\) is in \term{Jordan normal form} if it is a block diagonal matrix of the form
    \[
        A = \pmqty{
            J_{n_1}\pqty{\lambda_1} & & & & 0 \\
             & J_{n_2}\pqty{\lambda_2} & & & \\
             & & J_{n_3}\pqty{\lambda_3} & & \\
             & & & \ddots & \\
            0 & &  &  & J_{n_k}\pqty{\lambda_k}
        },
    \]
    where \(J_{n_i}\pqty{\lambda_i}\) are Jordan blocks of size \(n_i\) corresponding to eigenvalues \(\lambda_i\), not necessarily distinct. 
\end{definition}